\section{Research Plan and Methods}

\subsection{Data sources}

The project will use deeply phenotyped cohorts with plasma proteomics,
clinical diagnosis, cognitive testing, and neuroimaging where
available.

\subsection{Proteomics preprocessing}

Protein measurements will undergo rigorous quality control,
normalization, and batch-effect assessment. Missingness patterns will
be characterized explicitly, and model development will be embedded in
leakage-free pipelines.

\subsubsection{Quality control}

Sample- and protein-level exclusions will follow pre-specified
criteria. Outliers, technical artifacts, and panel-specific
inconsistencies will be documented.

\subsubsection{Normalization and harmonization}

We will compare within-platform normalization strategies and
cross-platform alignment approaches where relevant.

\subsection{Model development}

We will evaluate both supervised and representation-learning
approaches. Internal validation will use nested cross-validation or
split-sample designs with strict separation between preprocessing,
feature selection, and evaluation.

\subsection{Statistical analysis}

Performance will be quantified using discrimination, calibration,
uncertainty, and clinical operating characteristics. Subgroup
robustness will be evaluated by age, sex, disease stage, and cohort.
